\chapter{System Requirements}
\label{ch:requirements}

This chapter enumerates the functional and non-functional requirements
that were used to steer the design of \projname. Each requirement is
tagged with an identifier of the form \code{FR-}$n$ or \code{NFR-}$n$
so that later chapters can refer to it explicitly.

\section{Functional Requirements}
\label{sec:fr}

Table~\ref{tab:fr} lists the functional requirements. Each requirement
maps to one or more implemented modules and is verified in
Chapter~\ref{ch:experiments}.

\begin{table}[H]
\centering
\caption{Functional requirements and their implementation locations.}
\label{tab:fr}
\small
\begin{tabularx}{\linewidth}{@{}lXl@{}}
\toprule
\textbf{ID} & \textbf{Requirement} & \textbf{Implemented in} \\
\midrule
FR-01 & Accept CSV, TSV, and Parquet uploads and store them locally. & \file{core/workspace.py}, \file{app/streamlit\_app.py} \\
FR-02 & List all uploaded datasets in reverse-chronological order. & \file{core/workspace.py::list\_datasets} \\
FR-03 & Ingest a dataset and return shape, dtypes, and a preview. & \file{mcp\_servers/data.py::ingest\_dataset} \\
FR-04 & Validate a dataset against an auto-inferred pandera schema. & \file{mcp\_servers/data.py::validate\_schema\_with\_pandera} \\
FR-05 & Run a full EDA and produce three named plots. & \file{mcp\_servers/eda.py::run\_full\_eda} \\
FR-06 & Render a standalone correlation matrix with Pearson, Spearman, or Kendall. & \file{mcp\_servers/eda.py::render\_correlation\_matrix} \\
FR-07 & Train four candidate models in parallel with $k$-fold CV. & \file{mcp\_servers/modeling.py::run\_parallel\_bake\_off} \\
FR-08 & Run a time-bounded Optuna TPE hyperparameter sweep. & \file{mcp\_servers/modeling.py::trigger\_hyperparameter\_sweep} \\
FR-09 & Compute global SHAP feature importances for the champion model. & \file{mcp\_servers/explain.py::calculate\_shap\_values} \\
FR-10 & Produce a native feature-importance bar chart for tree models. & \file{mcp\_servers/explain.py::generate\_feature\_importance\_plot} \\
FR-11 & Emit a runnable Jupyter notebook that reproduces the pipeline. & \file{mcp\_servers/export.py::generate\_jupyter\_notebook} \\
FR-12 & Compile a PDF report of the champion, leaderboard, and metadata. & \file{mcp\_servers/export.py::compile\_pdf\_report} \\
FR-13 & Expose a unified tool catalog to the LLM and route calls to the right server. & \file{orchestrator/pool.py::MCPClientPool} \\
FR-14 & Stream tool progress and tool results back to the user in real time. & \file{orchestrator/loop.py::stream\_chat} \\
FR-15 & Preview and download every generated artefact. & \file{app/streamlit\_app.py} (\textit{Artifacts} tab) \\
FR-16 & Expose the orchestrator as a Python API for scripted use. & \file{orchestrator/\_\_init\_\_.py} \\
\bottomrule
\end{tabularx}
\end{table}

\section{Non-Functional Requirements}
\label{sec:nfr}

Table~\ref{tab:nfr} lists the non-functional requirements.

\begin{table}[H]
\centering
\caption{Non-functional requirements.}
\label{tab:nfr}
\small
\begin{tabularx}{\linewidth}{@{}lXX@{}}
\toprule
\textbf{ID} & \textbf{Requirement} & \textbf{Realisation} \\
\midrule
NFR-01 & Single-machine deployment; no cloud dependency for storage. & Local FS registry (\file{core/workspace.py}). \\
NFR-02 & No authentication or user model. & Removed in the migration; enforced by absence of any auth module. \\
NFR-03 & Fully open dependency stack. & All deps in \file{requirements.txt} are OSS. \\
NFR-04 & Deterministic numerical results under fixed \code{random\_state}. & \code{random\_state} arguments plumbed through all training tools. \\
NFR-05 & Progress feedback within 2\,s of tool invocation. & Progress bridge in \file{orchestrator/loop.py} + \code{ctx.report\_progress}. \\
NFR-06 & Extensible tool catalog. & Register-decorator MCP tools; no orchestrator changes for new tools. \\
NFR-07 & Reproducible artefacts. & UUID + kind + metadata in \file{data/artifacts/\_index.json}. \\
NFR-08 & Streamlit UI free of business logic. & UI only calls \code{orchestrator.stream\_chat} and \file{core/workspace.py}. \\
NFR-09 & Configurable data directory. & \code{DATA\_DIR} env variable via \code{pydantic-settings}. \\
NFR-10 & Upload size cap to prevent OOM. & \code{MAX\_UPLOAD\_MB} setting; enforced client-side and in tools. \\
NFR-11 & Reasonable page load / responsiveness on a laptop. & Streamlit + in-process MCP eliminates network round-trips. \\
NFR-12 & Ability to run without an LLM key (tool catalog still visible). & Streamlit degrades gracefully when \code{GROQ\_API\_KEY} is missing. \\
\bottomrule
\end{tabularx}
\end{table}

\section{Environmental Requirements}
\label{sec:env-req}

\begin{itemize}[leftmargin=1.4em]
  \item \textbf{Operating system:} Any OS with Python 3.11+; the
        development environment was Windows 11 with PowerShell, but
        no Windows-only APIs are used.
  \item \textbf{Python:} 3.11 or newer. Older versions lack the
        \code{X | Y} union syntax used throughout the codebase.
  \item \textbf{Memory:} 8~GB RAM is comfortable for datasets under
        the default 200~MB upload cap; SHAP on wide datasets can push
        this higher.
  \item \textbf{Network:} Outbound HTTPS to \code{api.groq.com} is
        required for the chat feature; no inbound network is used.
  \item \textbf{Browser:} Any modern browser for the Streamlit UI at
        \code{localhost:8501}.
\end{itemize}

\section{Traceability}
\label{sec:trace}

Table~\ref{tab:trace} traces each requirement to the chapter in which
its implementation is discussed. Chapter~\ref{ch:results} evaluates
the extent to which each requirement is satisfied.

\begin{table}[H]
\centering
\caption{Requirement--chapter traceability.}
\label{tab:trace}
\begin{tabular}{@{}lll@{}}
\toprule
\textbf{Requirement} & \textbf{Discussed in} & \textbf{Verified in} \\
\midrule
FR-01, FR-02 & Ch.~\ref{ch:design} & Ch.~\ref{ch:experiments} \\
FR-03, FR-04 & Ch.~\ref{ch:pipeline}, Ch.~\ref{ch:impl} & Ch.~\ref{ch:experiments} \\
FR-05, FR-06 & Ch.~\ref{ch:pipeline} & Ch.~\ref{ch:results} \\
FR-07, FR-08 & Ch.~\ref{ch:pipeline}, Ch.~\ref{ch:impl} & Ch.~\ref{ch:results} \\
FR-09, FR-10 & Ch.~\ref{ch:pipeline} & Ch.~\ref{ch:results} \\
FR-11, FR-12 & Ch.~\ref{ch:pipeline} & Ch.~\ref{ch:demo} \\
FR-13, FR-14 & Ch.~\ref{ch:mcp} & Ch.~\ref{ch:results} \\
FR-15, FR-16 & Ch.~\ref{ch:design}, App.~\ref{app:api} & Ch.~\ref{ch:demo} \\
NFR-01--NFR-12 & Ch.~\ref{ch:arch}, Ch.~\ref{ch:design} & Ch.~\ref{ch:results}, Ch.~\ref{ch:limitations} \\
\bottomrule
\end{tabular}
\end{table}
